\documentclass[12pt,a4paper]{report}

\usepackage[T1]{fontenc}
\usepackage[utf8]{inputenc}
\usepackage[english]{babel}
\usepackage{lmodern}


\usepackage{graphicx}
\usepackage{geometry}
\usepackage{setspace}
\usepackage{amsmath}
\usepackage{amssymb}
\usepackage{bm}
\usepackage[numbers]{natbib}
\usepackage{hyperref}
\usepackage{booktabs}
\usepackage{longtable}
\usepackage{array}
\usepackage{microtype}
\emergencystretch=1em

\usepackage{etoolbox}
\usepackage{titlesec}

\geometry{
  a4paper,
  left=25mm,
  right=25mm,
  top=25mm,
  bottom=25mm
}

\onehalfspacing

\makeatletter
\patchcmd{\chapter}{\clearpage}{\par}{}{}%
\patchcmd{\chapter}{\cleardoublepage}{\par}{}{}%
\makeatother

\titlespacing*{\chapter}{0pt}{1.5ex plus 1ex minus 0.5ex}{2ex plus 1ex}

\title{\textbf{Nowcasting of GDP: International Experience, Methodology, and Evaluation}\\
\large Master Thesis - Part I}
\author{}
\date{\today}

\begin{document}

\maketitle
\tableofcontents

%=========================================================
\chapter{Research Overview and Theoretical Foundations}

\section{Background, Motivation and Research Gap}

Among the central challenges faced by modern macroeconomic analysis is
the persistent lag with which official statistics measure the state of
real economic activity. Gross Domestic Product (GDP)---the most widely
used indicator of aggregate output---is reported quarterly and typically
published with a delay of four to eight weeks following the end of the
reference period. During this window, monetary authorities, fiscal
planners, and financial market participants must form judgements about
the current economic conjuncture on the basis of data that may already
be considerably outdated. This information deficit takes on particular
significance during episodes of rapid change, when waiting for official
releases can delay essential policy responses by an entire quarter or
more.

The concept of \textit{nowcasting}---predicting the present, the very
near future, and the very recent past of a target variable---has emerged
as a systematic response to this challenge. Originally a term from
meteorology, it was formally imported into economics by Giannone,
Reichlin, and Small, who demonstrated how the continuous flow of monthly
and higher-frequency statistical releases could be exploited to produce
timely estimates of quarterly GDP before the official figure becomes
available.\citep{Giannone2008} Unlike conventional medium-term
forecasting, nowcasting is explicitly oriented toward real-time data
assimilation: each new release of industrial production, survey indices,
or financial variables is treated as a signal that revises the current
assessment of aggregate output.

Although the nowcasting literature is mature and well-developed for
advanced economies---particularly the United States and the euro
area---a less-examined question concerns the relative performance of
traditional econometric frameworks versus modern machine learning (ML)
approaches, especially in emerging and transition economies where data
are scarcer, structural breaks are more frequent, and institutional
capacity differs substantially. The informational value of unstructured
digital data, such as internet search volumes, also remains inadequately
explored in economies where available indicators and internet
connectivity differ from those of high-income countries. This thesis
addresses these gaps through a comparative synthesis drawing on evidence
from six economies: the Netherlands, the United States, Peru, Georgia,
Armenia, and Belarus.

\section{Research Questions, Contributions and Analytical Scope}

Three core research questions organise the enquiry. First, how do
classical econometric models---bridge equations, mixed-data sampling
(MIDAS) regressions, and Dynamic Factor Models (DFMs)---compare to
machine learning algorithms in terms of point nowcast accuracy for GDP
at very short horizons? Second, do non-linear and high-dimensional
models such as random forests, support vector machines, and recurrent
neural networks confer particular advantages during episodes of
structural volatility, including the global financial crisis and the
COVID-19 pandemic? Third, can forecast combination strategies and model
interpretability tools---most notably Shapley value
decompositions---extend the practical utility of nowcasts beyond what
any single model delivers?\citep{NetherlandsDNB2022}

The thesis makes contributions at three levels.
\textbf{Methodologically}, it provides a systematic head-to-head
comparison of traditional and ML-based nowcasting specifications applied
within a unified pseudo-real-time evaluation framework.
\textbf{Empirically}, it documents model performance across distinct
phases of the business cycle, distinguishing the behaviour of competing
approaches in tranquil periods from their behaviour during crisis
episodes. \textbf{Practically}, it assesses whether ensemble
combinations and transparency tools can render complex ML nowcasts
sufficiently interpretable for day-to-day macroeconomic surveillance.

The analytical scope is restricted to short-horizon nowcasting of real
GDP growth for a single country case study, with the sample period
determined by the joint availability of quarterly national accounts and
a sufficiently rich set of monthly indicators. Three simplifying
assumptions apply throughout: a \textit{pseudo-real-time} design that
replicates real-time information flow using final-vintage data combined
with historical release calendars; approximate stationarity achieved
through standard transformations (logarithms, differencing, seasonal
adjustment); and a finite evaluation window that may encompass only a
limited number of crisis observations, constraining the statistical
power of state-dependent comparisons.

%=========================================================
\chapter{International Experience of Nowcasting}

\section{Nowcasting: Definition, Data and Modelling Strategies}

At its core, the nowcasting problem arises from a fundamental temporal
mismatch: the variable of primary interest---GDP---is released with a
lag of one to two months following the close of the quarter, whereas a
large and diverse set of related indicators becomes available on a
monthly, weekly, or even daily basis throughout that same quarter.
Giannone et al.\ formalise this as a problem of optimal filtering: the
nowcast of \(y_t\) is the best statistical estimate given the
information set \(\Omega_t\) available at calendar time \(t\), where
\(\Omega_t\) contains different series observed up to different points
in time---producing a dataset with a characteristic ``ragged edge'' at
its right-hand boundary.\citep{Giannone2008} The nowcaster's task is
therefore to process an asynchronous and growing information set in a
way that is both statistically coherent and computationally feasible in
real time.

A key analytical concept emerging from this framework is the notion of
\textit{news}: when a new data release arrives, its contribution to
revising the nowcast can be decomposed into two components.
\textit{Timeliness} captures the informational value derived simply from
arriving earlier than other indicators. \textit{Economic content}
captures the intrinsic explanatory power of the series for GDP once
the timing advantage is controlled for. Survey indicators---such as the
Federal Reserve Bank of Philadelphia's business outlook survey---carry a
disproportionately large impact on real-time nowcast revisions compared
with the Employment Report, precisely because they arrive earlier in
the publication cycle. Once timeliness is controlled for, however, hard
data like industrial production and retail sales prove highly significant
in terms of economic content. This distinction has important practical
implications: the optimal response to a new data release is not uniform
but depends on its location in the publication calendar.

Four broad families of indicators populate a typical nowcasting dataset.
\textbf{Hard macroeconomic data}---industrial production, retail and
wholesale trade, employment, exports and imports---are direct components
or close proxies of the national accounts, making them the most
informative signals for GDP once released. \textbf{Soft survey-based
indicators}---purchasing managers' indices, business tendency surveys,
consumer confidence---react more quickly to news but may embed
subjective noise. \textbf{Financial variables}---yield spreads, equity
indices, credit volumes, exchange rates---are available continuously and
encode forward-looking market assessments. \textbf{Alternative and
high-frequency data}---Google Trends search volumes, payment card
records, mobility indices---are especially relevant at the start of a
quarter when conventional statistics are absent.\citep{Peru2024,US2022}

Three modelling strategies dominate the literature for handling
mixed-frequency data. \textit{Bridge equations} aggregate monthly
predictors to quarterly frequency and estimate a direct projection of
GDP, providing a transparent but timing-limited approach.
\textit{MIDAS} regressions retain the original high-frequency
observations through parsimonious distributed-lag polynomial weights.
Most powerfully, \textit{Dynamic Factor Models} in state-space form use
the Kalman filter to handle the ragged edge directly, simultaneously
extracting latent common factors from a large panel and updating
nowcast estimates as each new data release arrives.\citep{Giannone2008}
Machine learning methods address the ragged edge by constructing
balanced design matrices through lagging or last-observation-carried-forward
imputation, making them applicable without modifying the core estimation
routine.

\section{International Experience: Advanced and Emerging Economies}

In advanced economies, central banks were among the first to
institutionalise systematic nowcasting. The Federal Reserve System and
the European Central Bank both deploy variants of the factor-model
framework, using panels of two hundred or more monthly series to produce
continuous updates of current-quarter GDP. Bok et al.\ provide a
comprehensive account of how such a system is operationalised in
practice at the Federal Reserve Bank of New York, maintaining a tracking
database of approximately 38{,}000 monthly and quarterly economic
series, automating the data ingestion pipeline, and formally quantifying
the ``news'' contribution of each scheduled release to the overall
nowcast revision.\citep{Bok2017} Their framework highlights that the
number of relevant latent factors need not be fixed in advance but can
be selected through information criteria applied to the full data panel,
a practical refinement over earlier implementations. A complementary
dimension of real-time analysis is documented by Anesti, Galv\~{a}o,
and Miranda-Agrippino, who show that the quality of preliminary nowcasts
depends partly on how substantially initial GDP releases are later
revised---implying that model rankings evaluated on final-vintage data
may look different from rankings obtained in genuine real
time.\citep{Anesti2020}

For the Netherlands, Kant, Pick and de Winter compare DFMs against a
suite of machine learning methods over 1992--2018, finding that the
Random Forest delivers the most accurate nowcasts for the current
quarter---an RMSFE of 0.91 in month~1 versus 0.98 for the DFM relative
to the prevailing mean benchmark---while the DFM recovers its advantage
for backcasting once hard data on the target quarter begin to arrive
(RMSFE 0.79 versus 0.88).\citep{NetherlandsDNB2022} Studies for Armenia
and Belarus demonstrate that LASSO (Armenia: RMSE~3.38) and SVM
(Belarus: RMSE~1.44) minimise forecast errors for real GDP, because
growth in both countries is governed by a narrow set of dominant
variables whose co-movement with GDP is captured efficiently by sparse
linear specifications; for the GDP deflator, Recurrent Neural Networks
achieve RMSE values of 2.16 and 1.05 respectively, because inflation
exhibits strong path-dependence and non-linear momentum that RNNs
capture through accumulated temporal context.\citep{ArmeniaBelarus2024}
In Peru, Gradient Boosting Machines incorporating Google Trends proxies
achieve an RMSE of 0.17---compared with 2.55 for the AR(1)
baseline---and successfully anticipated the output contraction of March
2020 weeks before official statistics confirmed it.\citep{Peru2024} In
Georgia, ML algorithms outperform the traditional DFM by 17--21~percent
in RMSE, with gains concentrated in periods incorporating VAT turnover
and Google Trends indices as high-frequency ``Type-1'' data available
within the first week of the following month.\citep{Georgia2022} In the
United States, a Bayesian Structural Time Series model augmented with
155 Google Trends search-volume categories reduces RMSFE by up to
40~percent in the first two months of the quarter---the period when
conventional hard indicators are almost entirely unavailable.\citep{US2022}

%=========================================================
\chapter{Why Nowcasting is Important}

\section{Policy, Market and Crisis Relevance}

The operational relevance of nowcasting derives principally from the
asymmetry between the pace at which economic conditions evolve and the
pace at which official measurement systems document those conditions.
For monetary authorities, real-time estimates of GDP growth and
inflation are essential inputs to interest-rate decisions and forward
guidance communications: because policy acts on the economy with a lag
of several quarters, the assessment of the \textit{current}
conjuncture---including the economy's position relative to potential
output---is at least as important as the medium-term forecast.
Nowcasting systems allow central banks to quantify how each individual
data release shifts the estimate of current-quarter GDP---the
\textit{news} released by each new data point---enabling a systematic
and reproducible real-time process in place of informal data
watching.\citep{Giannone2008}

A particularly important but underappreciated dimension of real-time
analysis is that policymakers face a \textit{double uncertainty}: they
must act not only on the basis of delayed data, but on the basis of data
that will itself be substantially revised in subsequent releases. Anesti
et al.\ document this for the United Kingdom, demonstrating that the
gap between preliminary and final-vintage GDP estimates is large enough
to affect the sign of quarterly growth in a non-trivial share of
observations, which in turn implies that real-time nowcast rankings
across models may look quite different from rankings obtained with
revised data.\citep{Anesti2020} This finding underscores why evaluating
models in a genuinely pseudo-real-time setting matters greatly for
drawing policy-relevant conclusions. Fiscal authorities face a parallel
need: revenue projections, budget balance estimates, and
debt-sustainability analyses all depend on the current level of economic
activity, and erroneous real-time estimates can lead to premature
tightening or unnecessarily delayed stimulus.

During the COVID-19 pandemic, nowcasting frameworks enabled governments
in the Netherlands, Georgia, and Armenia to calibrate emergency fiscal
packages to an estimate of the output decline that preceded official GDP
data by an entire quarter. Crucially, machine learning models proved
most valuable precisely during the most severe episodes: the Netherlands
study finds that the Random Forest's advantage over the DFM was most
pronounced both during the COVID-19 period and during the global
financial crisis of 2008--2009---precisely the episodes where the
non-linear, regime-switching behaviour of the economy fell outside the
range of experience that linear factor models had been trained to
capture.\citep{NetherlandsDNB2022} Financial market participants constitute
a second major constituency: the BSTS model of Kohns and Bhattacharjee
shows that incorporating Google Trends data produces materially sharper
estimates of US GDP in the weeks before any hard data are available,
allowing market participants to update yield and equity price
expectations with better information.\citep{US2022}

%=========================================================
\chapter{Analytical and Data Framework}

\section{Econometric and Machine Learning Methodology}

The Dynamic Factor Model occupies the foundational position in
central-bank nowcasting practice because it simultaneously addresses the
curse of dimensionality and the ragged-edge problem. In its canonical
formulation, the co-movement of a large \(N \times 1\) vector of
standardised macroeconomic series \(X_t\) is driven by a small
\(r \times 1\) vector of latent common factors \(F_t\) through the
\textbf{observation equation}
\begin{equation}
    X_t = \Lambda F_t + \varepsilon_t, \quad \varepsilon_t \sim \mathcal{N}(0,\Psi),
\end{equation}
where \(\Lambda\) is the \(N \times r\) matrix of factor loadings and
\(\varepsilon_t\) is the vector of idiosyncratic disturbances. The
factor dynamics follow the \textbf{transition equation}
\begin{equation}
    F_t = A F_{t-1} + u_t, \quad u_t \sim \mathcal{N}(0,Q),
\end{equation}
and the full system is estimated by the Kalman filter, which provides
the minimum mean-squared-error estimate of \(F_t\) conditional on
observed data, handling missing observations at the ragged edge by
skipping the measurement update for unreleased series.\citep{Giannone2008}
From an economic standpoint, the latent factors can frequently be mapped
to interpretable macro-financial concepts---a global real activity
factor, a financial conditions factor, a price-pressure factor---each
capturing a dimension of the business cycle as documented in the
large-scale operational implementation at the Federal Reserve Bank of
New York.\citep{Bok2017}

Machine learning approaches extend the nowcasting toolkit by relaxing
linearity and sparsity restrictions. \textit{Penalised
regressions}---LASSO \((\hat{\beta} = \arg\min_\beta \{\|y -
X\beta\|^2 + \lambda\|\beta\|_1\})\) and Ridge---shrink or zero out
coefficients of less informative predictors, performing automatic
variable selection or multicollinearity handling respectively.
\textit{Random Forests} grow an ensemble of \(B\) regression trees on
bootstrapped samples and random predictor subsets, averaging their
predictions \(\hat{f}_{RF}(x) = B^{-1}\sum_{b=1}^{B}\hat{f}_b(x)\)
to capture non-linear regime dependencies and interaction
effects.\citep{NetherlandsDNB2022} \textit{Recurrent Neural Networks}
maintain a hidden state \(h_t = \sigma(W_{hh}h_{t-1} + W_{xh}x_t)\)
that accumulates temporal information across time steps, making them
well-suited to persistent and non-linear price dynamics. When individual
model predictions are further combined with inverse-RMSE weights, the
resulting ensemble typically outperforms any single
specification---a pattern confirmed across the Armenia/Belarus and
Georgian country studies.\citep{ArmeniaBelarus2024}

\section{Data Infrastructure, Transformations and Feature Engineering}

A nowcasting dataset integrates indicators spanning multiple
frequencies. \textit{Real-activity hard data}---industrial production,
retail trade indices, exports, imports, employment---are the most
directly relevant because they measure components or close analogues of
the national expenditure accounts, but are subject to publication lags
of three to six weeks. \textit{Soft survey indicators}---PMIs, business
tendency and consumer confidence surveys---are released faster but embed
subjective reporting bias. \textit{Financial market variables}---term
spreads, equity indices, credit aggregates, exchange rates---are
available in real time and encode forward-looking expectations.
\textit{Alternative high-frequency data}---Google Trends indices,
card-transaction aggregates, mobility records---are particularly
valuable at the very start of a quarter before any conventional monthly
statistics are published.\citep{Georgia2022}

Preparing these heterogeneous series for modelling requires a structured
sequence of transformations. Level variables are converted to
log-differences or growth rates to achieve approximate stationarity;
seasonal and trading-day effects are removed using standard filter-based
decompositions; all series are then standardised (zero mean, unit
variance) before entry into factor models or penalised regressions.
Pseudo-real-time evaluation demands an additional layer of discipline: a
release calendar for each series must be reconstructed from historical
records so that each nowcast uses only data that would genuinely have
been available at the corresponding point in time.\citep{Anesti2020}
Feature engineering further enriches the predictor set: nominal trade
and retail series are deflated to real quantities; real interest rates
are computed by subtracting expected inflation; and Google Trends
indices are grouped into interpretable economic categories---labour-market
stress, consumption activity, financial distress---and de-seasonalised
before entry into the model.\citep{US2022}

%=========================================================
\chapter{Cross-Country Evidence and Comparative Assessment}

\section{Evidence from Developing Economies: Peru and Georgia}

Tenorio and Perez examine GDP nowcasting for Peru using a dataset
spanning January 2007 to May 2023, combining 91 structured
macroeconomic indicators with 10 groups of Google Trends variables
carefully designed to reflect distinct economic mechanisms: consumption
behaviour is proxied by brand-level queries (searches for ``Kia'',
``Toyota'', ``Restaurants'', ``Cinema''), financial conditions by
loan-market searches (``Credits'', ``Loans'', ``Mortgages'',
``Deals''), and external demand by travel-related and climate-event
queries (``flights'', ``visa'', ``El~Ni\~{n}o'').\citep{Peru2024}
Across a pseudo-real-time evaluation, Gradient Boosting Machines achieve
an RMSE of 0.17 compared with 0.32 for the best-performing DFM variant
and 2.55 for the AR(1) baseline---a reduction in forecast error of
more than 20--25~percent relative to conventional benchmarks. The
decisive advantage of the ML approach is most visible in two situations:
first, in the weeks immediately following a quarter's opening when hard
data remain entirely unavailable and only search-index signals carry
information; and second, during the COVID-19 shock of early 2020, when
models incorporating unstructured sentiment data correctly anticipated
the magnitude of the contraction weeks before official statistics
confirmed it. A key auxiliary finding is that combining structured and
unstructured predictors delivers gains exceeding what either data type
achieves in isolation, confirming a complementarity between conventional
macroeconomic indicators and digital footprint data.

Mgebrishvili's investigation of Georgia draws on a granular monthly
dataset that distinguishes ``Type-1'' indicators---exchange rates,
interbank rates, and tax revenues available within one week of
month-end---from ``Type-2'' indicators, including VAT turnover of
enterprises, tourism receipts, and remittances published around
mid-month; Google Trends search category indices are incorporated
alongside both types of conventional indicators as a separate
alternative data source.\citep{Georgia2022} Among the ensemble of ML
methods evaluated---LASSO, Ridge, Support Vector Regression, Linear
SVM, Neural Networks, XGBoost, CatBoost, AdaBoost, and Stacked
Generalisation---the ML ensemble collectively outperforms the DFM by
17--21~percent in RMSE on average, with the gap widening substantially
during the 2020--2021 pandemic episode. Whereas DFMs correctly identify
the sign of quarterly GDP growth in roughly two-thirds of quarters
under normal conditions, the best ML specifications achieve directional
accuracy above 70~percent even during turning-point episodes. Stacked
Generalisation---which combines the predictions of multiple base
learners through a meta-model trained on held-out data---proved
especially robust at turning points, suggesting that ensemble
aggregation across heterogeneous model families captures complementary
aspects of the economic signal.

\section{Evidence from Transition and Advanced Economies: Armenia, Belarus, Netherlands, and United States}

Tsukarev, Poghosyan, Lemba and Grishin conduct an exhaustive model
horse race for Armenia and Belarus, pitting bridge equations, MIDAS,
and DFM against Ridge, LASSO, Elastic Net, Bagging, Random Forest
(XGBoost), SVM, MLP, LSTM, and RNN.\citep{ArmeniaBelarus2024} One of
the study's headline findings is that machine learning algorithms
outperform conventional econometric models in every single tested
specification for both countries---a 100~percent win rate for ML across
all target variables and forecast horizons. For real GDP nowcasting,
LASSO attains the lowest RMSE in Armenia (3.38) and SVM dominates in
Belarus (1.44): growth cycles in both countries are governed by a narrow
set of drivers---commodity revenues, remittances, and sectoral
production---whose co-movement with GDP is captured efficiently by
sparse linear specifications that select these drivers while discarding
the many near-collinear indicators that add only noise. For the GDP
deflator, RNNs achieve RMSE values of 2.16 (Armenia) and 1.05
(Belarus), outperforming all competing methods, because inflation
exhibits highly path-dependent dynamics with non-linear momentum that
RNNs' recurrent hidden state captures through accumulated temporal
context. These contrasting results---sparse linear models for real
output, recurrent non-linear models for prices---illustrate a broader
principle: the optimal model architecture is determined not just by the
country context but by the specific statistical properties of the target
series. A combined nowcast averaging predictions across top-performing
ML models with inverse-RMSE weights provides the most stable
out-of-sample performance overall, confirming that diversification
across model families is a reliable hedge against individual
specification failures.

Kant, Pick and de Winter evaluate eleven methods for the Netherlands
over 1992Q1--2018Q4.\citep{NetherlandsDNB2022} At month-1 of the nowcast
horizon, Random Forest produces an RMSFE of 0.91 versus 0.98 for DFM.
The non-linearities captured by Random Forest are most valuable when
data are sparse: soft indicators (confidence surveys, PMIs) dominate
the predictor set early in the quarter and their relationship with GDP
changes non-linearly across business-cycle phases. Once the quarter
closes and hard data flow in, the DFM recovers its advantage, achieving
an RMSFE of 0.79 at the backcast horizon versus 0.88 for Random Forest.
Shapley value analysis reveals that variable importance shifts
systematically: survey-based soft indicators dominate SHAP contributions
during nowcasting, while industrial production and trade statistics
account for the majority of predictive weight at backcasting horizons.
For the United States, Kohns and Bhattacharjee's BSTS model augmented
with 155 Google Trends topic-level search indices---selected via the
Horseshoe shrinkage prior---reduces RMSFE by up to 40~percent in the
first two months of the quarter, with economy-related news, investing
behaviour, and labour-market proxy categories carrying the highest
posterior inclusion probabilities, confirming that internet-search
behaviour provides an almost-instantaneous proxy for GDP sentiment that
hard statistics will only document weeks later.\citep{US2022}

%=========================================================
\begin{thebibliography}{99}

\bibitem{Giannone2008}
Giannone, D., Reichlin, L., \& Small, D. (2006).
\href{https://www.ecb.europa.eu/pub/pdf/scpwps/ecbwp633.pdf}{\textit{Nowcasting: The real-time informational content of macroeconomic data}}.
European Central Bank Working Paper No.\ 633.

\bibitem{NetherlandsDNB2022}
Kant, D., Pick, A., \& de Winter, J. (2022).
\href{https://www.dnb.nl/media/kq4pe4cr/working_paper_no_754.pdf}{\textit{Nowcasting GDP using machine learning methods}}.
De Nederlandsche Bank Working Paper No.\ 754.

\bibitem{ArmeniaBelarus2024}
Tsukarev, T., Poghosyan, K., Lemba, K., \& Grishin, D. (2024).
\href{https://efsd.org/upload/iblock/0f4/hwg0io77bdcni8o7z0qmqjlpu7zagtpo/GDP-NOWCASTING-FROM-TRADITIONAL-ECONOMETRIC-MODELS-TO-MACHINE-LEARNING-ALGORITHMS.pdf}{\textit{GDP Nowcasting: From Traditional Econometric Models to Machine Learning Algorithms}}.
Eurasian Fund for Stabilization and Development, Working Paper WP/25/1.

\bibitem{Georgia2022}
Mgebrishvili, M. (2022).
\href{https://drive.google.com/file/d/118GZO60e9TBU52rk-vnPM0KlXjtRH4S2/view?usp=sharing}{\textit{Nowcasting with high frequency data, Google trends and machine learning}}.
Ministry of Finance of Georgia, Working Paper WP 01/2022.

\bibitem{Peru2024}
Tenorio, J., \& Perez, W. (2024).
\href{https://drive.google.com/file/d/1tIbCTI5vBiTK9OfBZ8C-3nT4Bnhuco1d/view?usp=sharing}{\textit{Monthly GDP nowcasting with Machine Learning and Unstructured Data}}.
arXiv:2402.04165v1.

\bibitem{US2022}
Kohns, D., \& Bhattacharjee, A. (2022).
\href{https://arxiv.org/abs/2011.00938}{\textit{Nowcasting Growth using Google Trends Data: A Bayesian Structural Time Series Model}}.
arXiv:2011.00938v2.

\bibitem{Bok2017}
Bok, B., Caratelli, D., Giannone, D., Sbordone, A., \& Tambalotti, A. (2017).
\href{https://www.newyorkfed.org/medialibrary/media/research/staff_reports/sr830.pdf}{\textit{Macroeconomic Nowcasting and Forecasting with Big Data}}.
Federal Reserve Bank of New York, Staff Report No.\ 830.

\bibitem{Anesti2020}
Anesti, N., Galv\~{a}o, A. B., \& Miranda-Agrippino, S. (2020).
\href{https://www.bankofengland.co.uk/-/media/boe/files/working-paper/2018/uncertain-kingdom-nowcasting-gdp-and-its-revisions.pdf}{\textit{Uncertain Kingdom: Nowcasting GDP and its Revisions}}.
Bank of England, Staff Working Paper No.\ 764.

\end{thebibliography}

\end{document}
